\chapter{Frontend}
\label{cap:frontend-teoria}

\begin{edital}
Tecnologias e práticas frontend web --- HTML, CSS, UX, Ajax, frameworks
(VueJS, Angular, React); padrões de frontend; SPA e PWA; UX (sistemas de
gestão de conteúdo, arquitetura de informação, portais, workflow,
acessibilidade, usabilidade).
\end{edital}

\section{HTML e CSS}

\begin{conceito}[Por que semântica HTML é assunto, não decoração]
Visualmente, \texttt{<div class="menu">} e \texttt{<nav>} podem
renderizar de forma idêntica --- a diferença é que só o segundo
\textbf{informa o papel} daquele bloco de conteúdo. Quem lê esse papel
são \textbf{programas}: o leitor de tela, que anuncia "navegação" e
permite ao usuário saltar direto para ela; o buscador, que pondera o
conteúdo de \texttt{<main>} de forma diferente do conteúdo de
\texttt{<footer>}; e o próprio navegador, que já entrega comportamento de
teclado pronto em elementos nativos (como \texttt{<button>} responder a
Enter e Espaço automaticamente). É daí que vem a "regra número um do
ARIA" vista adiante: se existe tag nativa com a semântica certa, ela
vence qualquer atributo acrescentado depois. A separação clássica ---
\textbf{HTML estrutura, CSS apresenta, JS comporta} --- é a mesma lógica
de baixo acoplamento aplicada à construção de uma página.
\end{conceito}

\begin{conceito}[Box model]
\textbf{content} $\to$ \textbf{padding} (espaço interno, dentro da
borda) $\to$ \textbf{border} (a borda em si) $\to$ \textbf{margin}
(espaço externo, fora da borda). O tamanho final \textbf{visível} de um
elemento é \texttt{content + padding + border} --- a \texttt{margin} conta
apenas para o espaço \textbf{ao redor} do elemento, não para o tamanho do
próprio elemento.

Com \texttt{box-sizing: border-box}, a propriedade \texttt{width} já
\textbf{inclui} padding e border dentro do valor declarado. Com o padrão
(\texttt{content-box}), \texttt{width} se refere só ao conteúdo, e
padding/border se \textbf{somam por fora} do valor declarado --- essa
diferença de \texttt{box-sizing} é onde nasce boa parte das pegadinhas de
cálculo de tamanho de elemento.
\end{conceito}

\begin{conceito}[Flexbox x Grid]
\textbf{Flexbox} é \textbf{unidimensional} --- distribui itens ao longo
de \textbf{um} eixo por vez (linha \textbf{ou} coluna). \textbf{Grid} é
\textbf{bidimensional} --- define linhas \textbf{e} colunas
simultaneamente. Um não substitui o outro: Grid é natural para a
\textbf{estrutura geral} da página; Flexbox é natural para
\textbf{alinhar o conteúdo} dentro de cada área já definida.
\end{conceito}

\begin{conceito}[Especificidade de seletores CSS]
Quando duas regras declaram a mesma propriedade para o mesmo elemento,
vence a de \textbf{maior especificidade}: \texttt{id} (peso 100) $>$
classe / atributo / pseudoclasse (peso 10) $>$ elemento (peso 1). A
\textbf{ordem de declaração} no arquivo só desempata quando as
especificidades são \textbf{iguais} --- uma regra de especificidade menor
declarada \textbf{depois} ainda perde para uma de especificidade maior
declarada antes. \texttt{!important} ignora toda essa hierarquia e vence
por cima de tudo, o que o torna desaconselhado como prática normal.
\end{conceito}

\begin{conceito}[\texttt{@import} x \texttt{<link>}]
Os dois trazem uma folha de estilo externa, mas \texttt{@import} é
resolvido \textbf{dentro do CSS}, em série (bloqueia o carregamento
seguinte até resolver), enquanto \texttt{<link>} no HTML permite
\textbf{download em paralelo} de várias folhas --- por isso
\texttt{<link>} é a prática recomendada em produção. \texttt{@import}
\textbf{não é} um mecanismo de responsividade, apesar de aceitar
condição de mídia como parâmetro.
\end{conceito}

\begin{cuidado}
Padding é interno, margin é externo --- não inverter. "A última regra
declarada sempre vence" é falso: só vence entre especificidades
\textbf{iguais}; um seletor de \texttt{id} ganha de um de classe mesmo
declarado antes dele no arquivo.
\end{cuidado}

\section{SPA x PWA}

\begin{conceito}
\begin{center}
\begin{tabular}{@{}p{2.6cm}p{4.6cm}p{4.6cm}@{}}
\toprule
& \textbf{SPA (\textit{Single Page Application})} & \textbf{PWA
(\textit{Progressive Web App})} \\
\midrule
O que é & app numa única página; troca conteúdo \textbf{sem recarregar} &
site que se comporta como um app nativo \\
Chave & roteamento no cliente, sem \textit{full reload} & \textbf{Service
Worker}, funcionamento offline, \textbf{instalável} \\
Depende de & JavaScript no navegador & recursos web progressivos \\
\bottomrule
\end{tabular}
\end{center}

O \textbf{Service Worker} é um script que roda \textbf{em segundo
plano}, separado da página em si, funcionando como um \textbf{proxy}
entre a aplicação e a rede --- é daí que vêm o cache e o funcionamento
offline: ele intercepta as requisições e decide servir da rede ou de um
cache local. Não tem acesso direto ao DOM e \textbf{exige HTTPS} (só
\texttt{localhost} é exceção durante desenvolvimento), precisamente
porque interceptar requisições numa conexão não cifrada seria um vetor
de ataque grave. O \textbf{manifest} (\texttt{manifest.json}) é o outro
pilar da PWA: define nome, ícones e modo de exibição que permitem a
\textbf{instalação} no dispositivo.
\end{conceito}

\begin{cuidado}
Service Worker é característica da \textbf{PWA}, não requisito de SPA
--- atribuí-lo à SPA é um erro comum. "SPA se instala no sistema
operacional como um app nativo" é também característica de \textbf{PWA},
não de SPA em si. Nenhuma das duas depende exclusivamente de um
framework específico --- ambas podem ser feitas com JavaScript puro.
\end{cuidado}

\subsection{Onde o HTML é montado: CSR, SSR, SSG e \textit{hydration}}

\begin{conceito}[O problema que a SPA criou, e as três respostas]
A SPA resolveu a navegação fluida, mas criou um problema novo: se o HTML
é montado inteiramente pelo JavaScript no navegador, o servidor entrega
uma página \textbf{vazia} de início, e o usuário espera o pacote de
código baixar e executar antes de ver qualquer conteúdo. Daí três
estratégias, distinguidas por \textbf{onde} e \textbf{quando} o HTML é
efetivamente gerado:

\begin{center}
\begin{tabular}{@{}p{3.1cm}p{3.5cm}p{4.6cm}@{}}
\toprule
\textbf{Estratégia} & \textbf{HTML é gerado\ldots} & \textbf{O que se
ganha e o que se paga} \\
\midrule
\textbf{CSR} (\textit{client-side}) & no \textbf{navegador}, a cada tela
& servidor leve e navegação instantânea depois da carga inicial;
\textbf{primeira} pintura lenta e indexação por buscadores mais frágil \\
\textbf{SSR} (\textit{server-side}) & no \textbf{servidor}, a cada
requisição & conteúdo visível e indexável de imediato; custo de CPU do
servidor a cada requisição \\
\textbf{SSG} (estático) & no momento do \textbf{\textit{build}}, uma
única vez & o mais rápido e barato de servir; o conteúdo só muda com uma
nova publicação \\
\bottomrule
\end{tabular}
\end{center}

\textbf{\textit{Hydration}} é o que costura SSR e SPA no mesmo produto: o
servidor manda HTML já pintado visualmente, mas HTML puro não tem
comportamento interativo --- quando o JavaScript chega ao navegador, ele
"hidrata" aquela marcação, reassociando componentes, estado interno e
ouvintes de evento a ela. A consequência prática que costuma virar item
de prova: existe uma \textbf{janela de tempo} em que a página
\textbf{parece pronta visualmente e não responde a cliques}, porque foi
pintada mas ainda não foi hidratada.

SSR/CSR decidem \textbf{quem monta o HTML}; Service Worker e cache
offline decidem \textbf{o que acontece quando não há rede} --- são eixos
\textbf{independentes}: uma SPA renderizada no cliente pode ser uma PWA
completa, e uma página com SSR pode não ser PWA nenhuma.
\end{conceito}

\section{Frameworks}

\begin{conceito}
\begin{center}
\begin{tabular}{@{}p{2.2cm}p{3.4cm}p{3.4cm}p{2.6cm}@{}}
\toprule
& \textbf{React} & \textbf{Angular} & \textbf{Vue} \\
\midrule
Tipo & biblioteca de UI (Meta) & framework completo (Google) & framework
progressivo \\
Linguagem & JS/JSX & TypeScript & JS \\
\bottomrule
\end{tabular}
\end{center}
React é deliberadamente uma \textbf{biblioteca focada em UI} (a
composição de outras peças --- roteamento, gerenciamento de estado global
--- fica a cargo do desenvolvedor); Angular é um \textbf{framework
completo}, com roteamento, injeção de dependência e formulários já
embutidos e opinativos sobre como o código deve ser organizado.
\end{conceito}

\begin{cuidado}
React não é escrito nativamente em TypeScript (embora aceite ser usado
com ele) --- o par React/framework-completo e Angular/biblioteca costuma
ser invertido em prova.
\end{cuidado}

\section{Ajax}

\begin{conceito}
\textbf{Ajax} (\textit{Asynchronous JavaScript and XML}, nome histórico)
é o conjunto de técnicas para fazer requisições \textbf{assíncronas} ao
servidor sem recarregar a página inteira --- é a base do comportamento de
uma SPA. Hoje implementado tipicamente via \texttt{fetch} ou
\texttt{XMLHttpRequest} (XHR), com dados trafegando majoritariamente em
JSON, apesar do "X" histórico do nome se referir a XML.
\end{conceito}

\section{UX, usabilidade e acessibilidade}

\begin{conceito}[Três recortes que se cruzam, mas não são sinônimos]
\begin{center}
\begin{tabular}{@{}p{2.9cm}p{4cm}p{4.5cm}@{}}
\toprule
\textbf{Recorte} & \textbf{Pergunta que responde} & \textbf{Como se
mede} \\
\midrule
\textbf{Usabilidade} & o usuário \textbf{típico} consegue cumprir a
tarefa com facilidade? & eficácia, eficiência e satisfação \textbf{num
contexto de uso} declarado \\
\textbf{Acessibilidade} & \textbf{pessoas com deficiência} conseguem
usar o sistema? & conformidade com critérios objetivos (WCAG A/AA/AAA) \\
\textbf{UX} & como é a \textbf{experiência inteira}, inclusive antes e
depois do uso propriamente dito? & percepção, expectativa, confiança ---
mais ampla e menos diretamente mensurável \\
\bottomrule
\end{tabular}
\end{center}

Duas relações resolvem a maior parte das questões sobre o tema.
\textbf{UX contém} usabilidade como um de seus atributos --- não são a
mesma coisa, e \textbf{UI} é apenas a camada visual através da qual a
experiência acontece, um subconjunto ainda menor. \textbf{Acessibilidade
não é um caso particular de usabilidade}: um sistema pode ser cômodo
para um usuário sem deficiência e \textbf{impossível} de usar para quem
navega por leitor de tela; e um sistema formalmente acessível (que passa
em todos os critérios técnicos) ainda pode ser péssimo de usar na
prática. Os dois conceitos se cruzam, não se contêm um no outro.

A diferença de \textbf{natureza}: usabilidade é objetivo de
\textbf{qualidade} de produto; acessibilidade é \textbf{requisito com
base normativa} --- WCAG (W3C) como critério técnico internacional,
\textbf{eMAG} como o modelo derivado dele adotado pelo governo
brasileiro, e a Lei Brasileira de Inclusão como obrigação legal em
serviços públicos.
\end{conceito}

\begin{conceito}[Nielsen x ISO 9241-11: duas listas, não misture]
Os \textbf{cinco atributos de usabilidade de Nielsen}: \textbf{facilidade
de aprendizado} (\textit{learnability}), \textbf{eficiência},
\textbf{memorabilidade} (o usuário que volta depois de um tempo relembra
como usar), \textbf{poucos erros} (e recuperáveis quando acontecem), e
\textbf{satisfação}. A \textbf{ISO 9241-11} usa outro trio, mais
compacto: \textbf{eficácia}, \textbf{eficiência} e \textbf{satisfação},
sempre avaliados \textbf{num contexto de uso} explicitamente declarado
(usuários, tarefas, equipamentos e ambiente específicos).
\end{conceito}

\begin{cuidado}
"Eficácia" pertence à lista da ISO 9241-11; "memorabilidade" e "poucos
erros" pertencem à lista de Nielsen --- misturar as duas listas é o erro
mais comum ao responder sobre o tema.
\end{cuidado}

\begin{conceito}[Acessibilidade: WCAG, níveis de conformidade, ARIA]
\textbf{WCAG} (W3C) organiza a acessibilidade em quatro princípios ---
\textbf{POUR}: \textbf{P}erceptível, \textbf{O}perável,
\textbf{C}ompreensível (\textit{Understandable}), \textbf{R}obusto. Cada
critério de sucesso individual tem um \textbf{nível de conformidade}:
\textbf{A} (mínimo), \textbf{AA} (o exigido na maioria das normas e
contratos públicos), \textbf{AAA} (máximo, nem sempre alcançável no site
inteiro sem comprometer outros objetivos de design). No Brasil, o
\textbf{eMAG} adapta esses critérios para o governo.

\textbf{ARIA} (\textit{Accessible Rich Internet Applications}): conjunto
de atributos (\texttt{role}, \texttt{aria-label}, \texttt{aria-hidden},
\texttt{aria-live}) que acrescenta \textbf{semântica} a elementos
\textbf{dinâmicos} ou construídos sem uma tag nativa correspondente, para
que leitores de tela anunciem função, nome e estado corretamente. A
\textbf{regra número um do ARIA é não usar ARIA}: se existe um elemento
HTML nativo com a semântica certa (\texttt{<button>}, \texttt{<nav>},
\texttt{<label>}), o nativo deve ser usado --- ARIA é complemento para o
que o HTML puro não cobre, e ARIA mal aplicado pode piorar a
acessibilidade em vez de melhorá-la (por exemplo, dizer a um leitor de
tela que algo é um botão quando ele não se comporta como um).
\end{conceito}

\begin{conceito}[Arquitetura de informação e CMS]
\textbf{Arquitetura de informação}: como o conteúdo de um site é
organizado, categorizado e navegado, de forma que o usuário encontre o
que procura. \textbf{CMS} (sistema de gestão de conteúdo): plataforma que
permite criar e gerenciar conteúdo sem escrever código a cada publicação
(WordPress, portais corporativos com fluxo editorial/\textit{workflow}
de aprovação).
\end{conceito}

\section{Mobile: multiplataforma x nativo}

\begin{conceito}
\textbf{Flutter} usa \textbf{Dart}; \textbf{React Native} usa
JavaScript; \textbf{Ionic} usa tecnologias web; \textbf{Xamarin} usa C\#.
\textbf{Mobile nativo}: Android usa \textbf{Kotlin} (linguagem
recomendada atualmente, com Java como opção mais antiga); iOS usa
\textbf{Swift} (com Objective-C como opção mais antiga).
\end{conceito}

\section{Leitura ativa de código: CSS e React}
\label{sec:leitura-ativa-frontend-teoria}

Boa parte das questões de frontend não pede definição --- pede para
\textbf{prever a renderização} de um trecho de HTML/CSS/JS. O protocolo
de leitura ativa do Capítulo \ref{cap:programacao-teoria} se aplica aqui:
resolva no papel, propriedade por propriedade, antes de comparar com
alternativas prontas.

\subsection{Pegadinhas visuais em CSS}

\begin{conceito}[A ordem certa para somar regras de CSS]
\begin{enumerate}
\item \textbf{Box model primeiro}: já visto acima --- confira
\texttt{box-sizing} antes de calcular qualquer tamanho.
\item \textbf{Pseudo-elementos}: \texttt{::before} insere conteúdo
\textbf{antes} do conteúdo real do elemento (mas ainda \textbf{dentro}
dele, no DOM renderizado); \texttt{::after}, \textbf{depois}. Sem uma
declaração \texttt{content: ...}, o pseudo-elemento \textbf{não aparece}
de jeito nenhum, mesmo que outras propriedades estejam definidas nele.
\item \textbf{\texttt{content: attr(x)}}: imprime o \textbf{valor} do
atributo \texttt{x} do elemento HTML como texto --- se o elemento não
tiver esse atributo, o pseudo-elemento aparece \textbf{vazio}, sem gerar
erro.
\item \textbf{Flex/Grid por último}: \texttt{flex-direction:
row-reverse} inverte a \textbf{ordem visual} dos itens, sem alterar a
ordem deles no DOM/HTML subjacente; \texttt{justify-content} age no
\textbf{eixo principal}, \texttt{align-items} no \textbf{eixo
transversal} --- trocar os dois é o distrator mais comum de layout.
\end{enumerate}
\end{conceito}

\begin{exemplo}
Um \texttt{<span data-nivel="urgente">} com CSS
\texttt{span::before \{ content: attr(data-nivel) ": "; font-weight:
bold; \}} produz, \textbf{antes} do texto original do \texttt{span}, o
literal \textbf{"urgente: "} em negrito --- não o texto
\texttt{"data-nivel"} (o nome do atributo, e não seu valor, que é o erro
mais comum de leitura de \texttt{attr()}).
\end{exemplo}

\begin{cuidado}
Questões de CSS costumam somar duas ou três regras na mesma pergunta
(\texttt{box-sizing} + pseudo-elemento + \texttt{flex-direction}, por
exemplo), e o erro de interpretação costuma estar na \textbf{combinação},
não em nenhuma regra isolada. Resolva cada camada separadamente antes de
somar o resultado final.
\end{cuidado}

\subsection{Escopo e \textit{closures} em JavaScript}

\begin{conceito}[O clássico laço com \texttt{var}]
\texttt{var} tem escopo de \textbf{função} (não de bloco) e sofre
\textit{hoisting} (é içada para o topo da função em que está). Dentro de
um laço, isso significa que existe \textbf{uma única} variável,
compartilhada por todas as funções criadas dentro do laço --- quando
essas funções forem finalmente executadas (por exemplo, num
\textit{callback} assíncrono), todas leem o valor \textbf{final} que a
variável tinha ao término do laço.

\begin{center}
\begin{tabular}{@{}p{5.4cm}p{5.4cm}@{}}
\toprule
\textbf{Código} & \textbf{Saída} \\
\midrule
\texttt{for (var i = 0; i < 3; i++)} \newline \texttt{\ \
fs.push(() => console.log(i));} & \textbf{3, 3, 3} --- um único \texttt{i}
compartilhado, que vale 3 ao final do laço \\
\texttt{for (let i = 0; i < 3; i++)} \newline \texttt{\ \
fs.push(() => console.log(i));} & \textbf{0, 1, 2} --- \texttt{let} cria
uma variável \textbf{nova a cada iteração} \\
\bottomrule
\end{tabular}
\end{center}
\end{conceito}

\begin{conceito}
\texttt{let} e \texttt{const} (ES6) têm escopo de \textbf{bloco}.
\texttt{const} impede a \textbf{reatribuição} da própria variável, mas
\textbf{não congela} o objeto referenciado --- \texttt{const obj = \{\};
obj.x = 1;} é perfeitamente válido, porque o que está protegido é a
ligação nome-valor, não o conteúdo interno do objeto.
\texttt{==} faz \textbf{coerção de tipo} antes de comparar
(\texttt{"1" == 1} resulta \texttt{true}); \texttt{===} compara valor
\textbf{e} tipo, sem coerção.
\end{conceito}

\begin{cuidado}
Um laço com \texttt{var} \textbf{não} imprime 0, 1, 2 --- essa saída
descreve o comportamento de \texttt{let}, não de \texttt{var}. Em Java, o
escopo de bloco é sempre a regra (Capítulo \ref{cap:java-teoria}) --- não
confundir o comportamento das duas linguagens.
\end{cuidado}

\subsection{Renderização no React}

\begin{conceito}
\textbf{\texttt{key} em listas}: React usa a \texttt{key} para
identificar qual item de uma lista mudou, foi adicionado ou removido
entre duas renderizações sucessivas. Usar o \textbf{índice do array}
como \texttt{key} numa lista que pode reordenar, inserir ou remover
itens \textbf{quebra} esse rastreamento (o React pode reaproveitar o
componente errado para o item errado, causando bugs visuais sutis) --- a
prática recomendada é usar um \textbf{identificador estável} vindo dos
próprios dados, nunca o índice posicional.

\textbf{\texttt{useState} é assíncrono e agrupado (\textit{batched})}:
chamar a função que atualiza o estado \textbf{não} atualiza a variável
\textbf{na mesma linha} de código --- o novo valor só se reflete na
\textbf{próxima renderização} do componente. Um trecho que chama
\texttt{setContador(contador + 1)} duas vezes seguidas \textbf{não} soma
2 ao contador, porque as duas chamadas leem o mesmo valor "antigo" de
\texttt{contador} dentro do mesmo ciclo de evento --- para acumular
corretamente usa-se a forma funcional \texttt{setContador(c => c + 1)},
que recebe o valor mais atualizado em cada chamada.

\textbf{Reconciliação (Virtual DOM)}: o React compara a árvore virtual
nova com a anterior e aplica no DOM real \textbf{só as diferenças} ---
não recria a página inteira a cada renderização. Criar uma
\textbf{função anônima nova a cada renderização} (por exemplo,
\texttt{onClick={() => ...}} sem memoização) não quebra a reconciliação
em si, mas pode gerar \textbf{re-renderizações desnecessárias} em
componentes filhos que foram otimizados com \texttt{memo}.
\end{conceito}

\begin{cuidado}
Usar o índice do array como \texttt{key} \textbf{não} é sempre seguro ---
falha justamente quando a lista pode reordenar ou ter itens
inseridos/removidos no meio. \texttt{setState}/\texttt{useState}
\textbf{não} atualiza de forma síncrona e imediata --- é assíncrono e
agrupado dentro do mesmo \textit{handler} de evento. Virtual DOM
(estrutura em memória usada para calcular diferenças) não é o mesmo que
o DOM real do navegador.
\end{cuidado}
